\clearpage
\VerdierPhysicalPageAnchor{0181}%
\begingroup\small\setlength{\parskip}{0pt}\setlength{\emergencystretch}{1.5em}
\setlength{\abovedisplayskip}{4pt plus 2pt minus 2pt}%
\setlength{\belowdisplayskip}{4pt plus 2pt minus 2pt}%

\VerdierSemanticAnchor{ch3d-1-2-10-proof-cont}
于是复形 $X(0,+\infty)$（(1.2.1) 的证明）与 $X$ 拟同构，并且典范态射
$H^0(X(0,+\infty))\longrightarrow X(0,+\infty)$ 是拟同构。这表明 $X$ 在
$D(\mathcal A)$ 中同构于 $H^0(X)$。

证明 b)。第二个等式列在这里是为了备忘，它由 $\mathcal A$ 是
$D(\mathcal A)$ 的满子范畴这一事实得出。$D(\mathcal A)$ 中从 $X[-n]$ 到
$Y$ 的一个态射具有 $Q(m)Q(s)^{-1}$ 的形式，其中
$m\colon Z^\bullet\longrightarrow Y$ 是 $K(\mathcal A)$ 中的态射，而
$s\colon Z^\bullet\longrightarrow X[-n]$ 是 $K(\mathcal A)$ 中的拟同构。
这时态射 $i\colon Z^\bullet(-\infty,n+1)\longrightarrow Z^\bullet$
（(1.2.1) 的证明）是拟同构。因此
\[
Q(m)Q(s)^{-1}=Q(m)Q(i)Q(i)^{-1}Q(s)^{-1}=Q(mi)Q(si)^{-1}.
\]
当 $n$ 严格小于零时，$mi=0$。因而 $Q(m)Q(s)^{-1}=0$。

\VerdierSemanticAnchor{ch3d-1-2-11}
\medskip\noindent\textbf{定义 1.2.11.}\quad
以 $\mathcal A^{\operatorname{Ext}}$ 表示 $D(\mathcal A)$ 的满子范畴，它由
$\mathcal A$ 的对象的各个平移定义。

因此，范畴 $\mathcal A^{\operatorname{Ext}}$（与子范畴 $\mathcal A$ 不同）
是 $D(\mathcal A)$ 的一个 $\mathbb Z$-子范畴。$D(\mathcal A)$ 的一个对象
同构于 $\mathcal A^{\operatorname{Ext}}$ 的一个对象，当且仅当它只有一个
非零上同调对象。

\VerdierSemanticAnchor{ch3d-1-3}
\subsection*{1.3. 典范 \texorpdfstring{$\delta$}{delta}-函子。}

\VerdierSemanticAnchor{ch3d-1-3-1}
\noindent\textbf{1.3.1.}\quad
设 $\mathcal A$ 是阿贝尔范畴，$\mathcal D$ 是三角范畴。从 $\mathcal A$
到 $\mathcal D$ 的一个 $\delta$-函子是由以下数据组成的对象：

\noindent 1) 一个函子 $G\colon\mathcal A\longrightarrow\mathcal D$。

\noindent 2) 对 $\mathcal A$ 中的每条正合列
$S=0\longrightarrow X\longrightarrow Y\longrightarrow Z\longrightarrow0$，
给定 $\mathcal D$ 中的一个态射：
\[
\delta(S)\colon G(Z)\longrightarrow G(X)[1].
\]
这些数据须满足以下公理：

\noindent a) 对每个正合列的态射：

\VerdierSemanticAnchor{ch3d-1-3-1-exact-sequences}
\begin{center}
\resizebox{0.88\textwidth}{!}{%
$\begin{tikzcd}[ampersand replacement=\&,column sep=2.6em,row sep=2.2em]
S={}\&0\arrow[r]\&X\arrow[r]\arrow[d,"u"']\&Y\arrow[r]\arrow[d,"v"']\&Z\arrow[r]\arrow[d,"w"']\&0\\
S'={}\&0\arrow[r]\&X'\arrow[r]\&Y'\arrow[r]\&Z'\arrow[r]\&0
\end{tikzcd}$%
}
\end{center}
\endgroup
